% Generated draft: requires mathematical and source review.
\documentclass[11pt]{article}
\usepackage[margin=25mm]{geometry}
\usepackage{fontspec}
\setmainfont{Latin Modern Roman}
\usepackage{amsmath,amssymb,mathtools,unicode-math}
\setmathfont{Latin Modern Math}
\usepackage{enumitem,needspace,xurl,hyperref}
\setlength{\parindent}{0pt}
\setlength{\parskip}{5pt}
\setlength{\emergencystretch}{3em}
\newcommand{\N}{\mathbb{N}}
\newcommand{\Z}{\mathbb{Z}}
\newcommand{\Q}{\mathbb{Q}}
\newcommand{\R}{\mathbb{R}}
\newcommand{\C}{\mathbb{C}}
\newcommand{\F}{\mathbb{F}}
\newcommand{\A}{\mathbb{A}}
\newcommand{\PP}{\mathbb{P}}
\newcommand{\E}{\mathbb{E}}
\newcommand{\eps}{\varepsilon}
\newcommand{\dd}{\,\mathrm{d}}
\begin{document}
\section*{UnsolvedMath ID: 700646}
\textbf{Source problem code:} catalogue entry
\subsection{Definitions and mathematical statement}
Let \(G\) be an algebraic group, let \(H\) be a closed normal subgroup of \(G\), and let
\[
Q=G/H,\qquad f:G\longrightarrow Q
\]
be the quotient and its canonical homomorphism. The supplied source does not specify the ground field or the precise category of algebraic groups and quotients; these choices remain unresolved.

For a fixed ground field \(k\) and an interpretation in which \(G\), \(H\), \(Q\), and \(f\) are defined over \(k\), a rational section of \(f\) means a rational map \(s:Q\dashrightarrow G\), defined over \(k\), represented by a morphism \(s_U:U\to G\) on a dense open subset \(U\subseteq Q\), such that
\[
f\circ s_U=\iota_U,
\]
where \(\iota_U:U\hookrightarrow Q\) is the inclusion. Equivalently, \(f\circ s=\operatorname{id}_Q\) as rational maps. This definition does not require \(s\) to be a group homomorphism or to be defined everywhere.

\textbf{Problem.} Determine necessary and sufficient conditions on \(G\), \(H\), and the relevant field of definition for the canonical homomorphism \(f:G\to G/H\) to admit a rational section.

\textbf{Unresolved source qualifications.} The source writes the section identity as \(sf=1\). The identity above uses the usual geometric meaning of a section; the source's order of composition is not explained. The ground field, whether the section must be defined over that field or after extending it, and any conventions concerning affineness, connectedness, smoothness, or reducedness are not supplied. No such extra hypotheses are imposed here; a definitive formulation requires these conventions to be fixed.
\subsection{Short English statement}
For an algebraic group \(G\) and a closed normal subgroup \(H\), determine when the quotient map \(G\to G/H\) has a rational section, meaning a map defined on a dense open subset that chooses a representative of each quotient point there. The ground field, field of definition, and composition convention require clarification.
\subsection{Sources}
\begin{itemize}
\item Local UnsolvedMath registry, permanent ID \(700646\), exact source problem code \texttt{catalogue entry}, title \emph{Kourovka Notebook Problem 10.31}. The supplied registry labels the problem \emph{open}; this label is not an independently verified mathematical or literature assessment.
\item \emph{The Kourovka Notebook (Unsolved Problems in Group Theory)}, original problem 10.31, 10th issue (1986), proposed by Yu. I. Merzlyakov. The supplied record identifies the reference edition as the 21st issue edition (2026), \texttt{arXiv:1401.0300v48}. The collection description names E. I. Khukhro and V. D. Mazurov as editors.
\item Yu. I. Merzlyakov, \emph{Rational Groups}, Nauka, Moscow, 1987, \(\S\,37\) (Russian). The supplied problem cites this section for partial results; their hypotheses and conclusions are not included in the record.
\end{itemize}
\end{document}
